\documentclass[11pt]{article}

\usepackage[utf8]{inputenc}
\usepackage[T1]{fontenc}
\usepackage{lmodern}
\usepackage{amsmath,amssymb,amsthm,amsfonts}
\usepackage{mathtools}
\usepackage{enumitem}
\usepackage{geometry}
\usepackage{hyperref}
\usepackage{mathrsfs}
\usepackage{microtype}

\geometry{margin=1in}
\numberwithin{equation}{section}

% --------------------------------------------------
% Theorem environments
% --------------------------------------------------
\newtheorem{theorem}{Theorem}[section]
\newtheorem{proposition}[theorem]{Proposition}
\newtheorem{lemma}[theorem]{Lemma}
\newtheorem{corollary}[theorem]{Corollary}

\theoremstyle{definition}
\newtheorem{definition}[theorem]{Definition}
\newtheorem{example}[theorem]{Example}

\theoremstyle{remark}
\newtheorem{remark}[theorem]{Remark}

% --------------------------------------------------
% Macros
% --------------------------------------------------
\newcommand{\Z}{\mathbb{Z}}
\newcommand{\R}{\mathbb{R}}
\newcommand{\C}{\mathbb{C}}
\newcommand{\N}{\mathbb{N}}

\newcommand{\dd}{\mathrm{d}}
\newcommand{\T}{\mathbb{T}}
\newcommand{\Id}{\mathrm{Id}}
\newcommand{\Hil}{\mathcal{H}}
\newcommand{\Htw}{\mathcal{H}_{\mathrm{twin}}}

\newcommand{\om}{\omega}
\newcommand{\Om}{\Omega}

\newcommand{\Gdisc}{G_{\mathrm{disc}}}
\newcommand{\Gcont}{G_{\mathrm{cont}}}

% --------------------------------------------------
% Title
% --------------------------------------------------
\title{\bf Foundations of Twin Quantum Geometry:\\
	Emergent Quantization from Discrete Orbital Structures}

\author{Jocelyn Nembe}
\date{}

\begin{document}
	
	\maketitle
	
	\begin{abstract}
		We develop a geometric foundation for quantum mechanics based on discrete
		group actions on phase space, formalized as twin spaces. Starting from a
		classical symplectic phase space $(M,\omega)$ and a discrete twin group
		$\Gdisc$ acting on $M$, we show that the Planck constant $\hbar$,
		canonical commutation relations, the Hilbert space structure and
		Schr\"odinger dynamics can be derived from orbital geometry rather than
		postulated. The Planck constant appears as the minimal non-trivial
		symplectic area of a fundamental twin cell, or equivalently as a
		cohomological invariant of the group action. Canonical commutation
		relations emerge from a non-trivial braiding of position and momentum
		translations on the twin lattice. We construct a natural twin Hilbert
		space $\Htw = L^2(M/\Gdisc,\mu)$ and prove that Schr\"odinger dynamics
		arises as covariant transport of Hamiltonian flows on the orbit space,
		under a quantization condition for the symplectic area. Several model
		systems (harmonic oscillator, free particle, particle on a circle)
		illustrate how standard quantum spectra are recovered within this purely
		geometric framework.
	\end{abstract}
	
	\bigskip
	
	\noindent\textbf{Keywords.}
	Twin spaces, discrete group actions, symplectic geometry,
	emergent quantization, canonical commutation relations,
	Hilbert space, Schr\"odinger equation.
	
	\medskip
	
	\noindent\textbf{MSC 2020.}
	81S10, 53D05, 81Q05, 37J05, 20H15.
	
	% ==================================================
	\section{Introduction}
	% ==================================================
	
	The standard axioms of quantum mechanics postulate the existence of
	\begin{itemize}[leftmargin=2em]
		\item a complex Hilbert space $\Hil$ of states;
		\item a distinguished constant $\hbar>0$;
		\item self-adjoint operators $Q,P,\dots$ satisfying canonical
		commutation relations, such as $[Q,P]=i\hbar \Id$;
		\item the Born rule, prescribing probabilities as squared norms of
		orthogonal projections.
	\end{itemize}
	In this picture, the Hilbert space, the value of $\hbar$, and the
	canonical structure are given \emph{a priori}. Even in geometric
	quantization and deformation quantization, quantization is encoded in
	additional structures rather than emerging from a more primitive
	geometry.
	
	In this article we propose a different viewpoint. We start from a
	classical symplectic phase space $(M,\omega)$ and a discrete group
	$\Gdisc$ acting on $M$ by symplectomorphisms, hence generating a
	lattice-like structure of orbits. We show that, under suitable
	assumptions, the following objects emerge from this orbital geometry:
	\begin{itemize}[leftmargin=2em]
		\item a distinguished area scale which we identify with $\hbar$;
		\item canonical commutation relations derived from the braiding of
		discrete position and momentum translations;
		\item a natural Hilbert space $\Htw=L^2(M/\Gdisc,\mu)$ associated with
		the orbit space and an invariant measure $\mu$;
		\item a unitary time evolution on $\Htw$ satisfying a Schr\"odinger-type
		equation arising from covariant transport of classical Hamiltonian
		flows.
	\end{itemize}
	
	This leads to the following main results.
	
	\begin{theorem}[Emergent Planck constant]\label{thm:intro-hbar}
		Let $(M,\omega)$ be a $2n$-dimensional symplectic manifold and let
		$\Gdisc$ be a discrete group of symplectomorphisms acting properly and
		cocompactly on $M$. Assume that the symplectic form has integral periods
		on fundamental homology cycles associated with the action. Then there
		exists a constant $\hbar>0$ such that the symplectic area of any
		fundamental cell of the action is an integer multiple of $2\pi\hbar$.
		The minimal non-zero area
		\[
		\Delta \mathcal{A}_{\min} = 2\pi \hbar
		\]
		is uniquely determined by the cohomology class $[\omega]$ and the action
		of $\Gdisc$.
	\end{theorem}
	
	\begin{theorem}[Geometric canonical commutation]\label{thm:intro-CCR}
		Let $T_q(a)$ and $T_p(b)$ denote discrete translations in position and
		momentum associated with generators of $\Gdisc$ along canonical
		directions in $M=\R^{2}$ or $M=\R^{2n}$. Under the quantization
		condition of Theorem~\ref{thm:intro-hbar} one has
		\[
		T_q(a)\, T_p(b) = e^{i ab/\hbar}\, T_p(b)\, T_q(a).
		\]
		In the infinitesimal limit $a,b\to 0$ this yields the canonical
		commutation relation
		\[
		[Q,P] = i\hbar \Id
		\]
		for densely defined self-adjoint generators $Q,P$ of these translations
		on a suitable Hilbert space.
	\end{theorem}
	
	\begin{theorem}[Twin Hilbert construction]\label{thm:intro-Hilbert}
		Assume that $\Gdisc$ acts properly discontinuously on $M$ and that the
		quotient $X:=M/\Gdisc$ carries a $\sigma$-finite measure $\mu$
		invariant under the induced dynamics. Then the space
		\[
		\Htw := L^2(X,\mu)
		\]
		is a separable Hilbert space, and the representations of translations,
		observables and Hamiltonian flows define unitary (or self-adjoint)
		operators on $\Htw$ compatible with the twin geometry.
	\end{theorem}
	
	\begin{theorem}[Emergent Schr\"odinger dynamics]\label{thm:intro-Schro}
		Let $H_{\mathrm{cl}}:M\to\R$ be a classical Hamiltonian whose flow
		preserves the quantized symplectic structure and commutes with the
		action of $\Gdisc$ up to gauge-equivalent phases. Then the induced
		evolution on $\Htw$ is unitary and satisfies a Schr\"odinger equation
		\[
		i\hbar \frac{\partial}{\partial t} \psi_t = \widehat{H}\, \psi_t,
		\]
		where $\widehat{H}$ is the twin-quantized Hamiltonian acting on
		$\Htw$.
	\end{theorem}
	
	The present article provides a detailed construction and proof of these
	results, together with illustrative examples. It serves as the
	geometric and analytic foundation for subsequent works:
	\begin{itemize}[leftmargin=2em]
		\item in \emph{Twin Quantum Mechanics: Operators, Measurement and
			Entanglement}, we develop the full operator, measurement and
		entanglement theory on $\Htw$;
		\item in \emph{Twin Quantum Fields, Kernel Dynamics and Gravity-like
			Effects}, we extend the construction to quantum fields, kernel
		dynamics and emergent gravity-like corrections.
	\end{itemize}
	
	\subsection*{Structure of the paper}
	
	In Section~\ref{sec:twin-phase} we introduce twin phase spaces and
	discrete group actions. Section~\ref{sec:hbar} derives the Planck
	constant from orbital geometry. In Section~\ref{sec:CCR} we construct
	geometric canonical commutation relations from the braiding of position
	and momentum translations. Section~\ref{sec:Hilbert} defines the twin
	Hilbert space as $L^2$ on the orbit space. Section~\ref{sec:Schro}
	derives Schr\"odinger dynamics from covariant transport of Hamiltonian
	flows. Finally, Section~\ref{sec:examples} presents simple models.
	
	
	% ==================================================
	\section{Twin phase spaces and discrete group actions}
	\label{sec:twin-phase}
	% ==================================================
	
	We start from a classical phase space and introduce its twin
	structure via a discrete group of symplectomorphisms.
	
	\subsection{Classical phase space}
	
	Let $(M,\omega)$ be a $2n$-dimensional symplectic manifold, where
	$\omega$ is a closed non-degenerate 2-form. Typical examples include:
	\begin{itemize}[leftmargin=2em]
		\item $M=\R^{2n}$ with canonical coordinates $(q,p)$ and
		$\omega = \sum_{j=1}^n \dd p_j \wedge \dd q_j$;
		\item the cotangent bundle $M = T^*Q$ of a configuration manifold $Q$,
		with its canonical symplectic form;
		\item compact phase spaces such as $T^*\T^n$ or coadjoint orbits.
	\end{itemize}
	
	\subsection{Discrete twin groups and orbits}
	
	\begin{definition}[Discrete twin group]
		A \emph{discrete twin group} on $(M,\omega)$ is a countable group
		$\Gdisc$ together with an injective morphism
		\[
		\rho : \Gdisc \to \mathrm{Symp}(M,\omega),
		\]
		where $\mathrm{Symp}(M,\omega)$ denotes the group of symplectomorphisms
		of $(M,\omega)$, such that $\rho(\Gdisc)$ acts properly discontinuously
		on $M$.
	\end{definition}
	
	We write $g\cdot x$ for the action of $g\in\Gdisc$ on $x\in M$. For each
	$x\in M$, its orbit is
	\[
	\mathcal{O}_x := \{g\cdot x : g\in\Gdisc\}.
	\]
	
	\begin{definition}[Twin phase space]
		The pair $(M,\Gdisc)$, together with the symplectic form $\omega$, is
		called a \emph{twin phase space}. The orbit space
		\[
		X := M/\Gdisc
		\]
		is the \emph{twin orbit space}. Points in $X$ correspond to orbits, or
		equivalently to equivalence classes $[x]=\mathcal{O}_x$ under the
		action of $\Gdisc$.
	\end{definition}
	
	We typically assume that the action is free or has mild stabilizers so
	that $X$ carries a reasonable differential or measurable structure.
	
	\subsection{Fundamental domains and cells}
	
	Under suitable assumptions, one can choose a fundamental domain for the
	action.
	
	\begin{definition}[Fundamental domain]
		A measurable subset $F\subset M$ is a \emph{fundamental domain} for the
		action of $\Gdisc$ if
		\begin{itemize}[leftmargin=2em]
			\item $M = \bigsqcup_{g\in\Gdisc} g\cdot F$ up to measure zero;
			\item $g\cdot F$ and $h\cdot F$ are disjoint up to measure zero if
			$g\neq h$.
		\end{itemize}
	\end{definition}
	
	In the examples of interest, $F$ can be chosen to have finite symplectic
	volume
	\[
	\mathrm{Vol}_\omega(F) := \int_F \omega^{\wedge n} < \infty.
	\]
	
	\begin{remark}[Twin cells]
		We interpret $F$ (or its translates) as fundamental \emph{twin cells}
		in phase space. In dimension $2$, these are area cells in the $(q,p)$
		plane; in higher dimension they are $2n$-dimensional volume cells.
		The minimal non-zero symplectic volume of such cells will be related to
		the Planck constant.
	\end{remark}
	
	
	% ==================================================
	\section{Emergence of the Planck constant from orbital geometry}
	\label{sec:hbar}
	% ==================================================
	
	We now show how a distinguished area scale arises from the symplectic
	structure and the discrete action.
	
	\subsection{Integral symplectic periods}
	
	Consider the de Rham cohomology class $[\omega]\in H^2_{\mathrm{dR}}(M)$
	and the image of integral homology cycles of the quotient. Let
	\[
	\pi : M \to X=M/\Gdisc
	\]
	denote the quotient map. Closed 2-cycles $\Sigma$ in $M$ induced by
	fundamental loops of the action may carry non-zero symplectic area
	\[
	\int_\Sigma \omega.
	\]
	
	\begin{definition}[Quantized symplectic action]
		We say that the symplectic form is \emph{quantized with respect to
			$\Gdisc$} if there exists a constant $\hbar>0$ such that
		\[
		\int_\Sigma \omega \in 2\pi\hbar\,\Z
		\]
		for all 2-cycles $\Sigma$ generated by the action of $\Gdisc$.
	\end{definition}
	
	\begin{remark}
		This condition is reminiscent of the prequantization condition in
		geometric quantization, but here it is imposed relative to the discrete
		group action rather than to all cycles in $H_2(M,\Z)$.
	\end{remark}
	
	\subsection{Emergent minimal area}
	
	Under this assumption, we can define the minimal non-zero area.
	
	\begin{proposition}[Minimal symplectic area]
		Assume that the set of non-zero symplectic areas
		\[
		\mathcal{S} := \left\{\int_\Sigma \omega \; :\;
		\Sigma \text{ is a 2-cycle generated by }\Gdisc \right\}
		\]
		is a discrete subset of $\R$. Then there exists a minimal non-zero value
		$\Delta\mathcal{A}_{\min}>0$ in $\mathcal{S}$.
	\end{proposition}
	
	\begin{proof}
		If $\mathcal{S}$ is discrete and contains non-zero elements, there is a
		smallest positive element by the well-ordering of discrete subsets of
		$\R_{\ge 0}$.
	\end{proof}
	
	\begin{theorem}[Emergent Planck constant]\label{thm:hbar}
		Under the quantization and discreteness assumptions above, there exists
		$\hbar>0$ such that every fundamental 2-cycle $\Sigma$ generated by
		$\Gdisc$ satisfies
		\[
		\int_\Sigma \omega \in 2\pi\hbar\,\Z,
		\]
		and the minimal non-zero area is
		\[
		\Delta\mathcal{A}_{\min} = 2\pi\hbar.
		\]
		The constant $\hbar$ is uniquely determined by the cohomology class
		$[\omega]$ restricted to the subgroup generated by the action of
		$\Gdisc$.
	\end{theorem}
	
	\begin{proof}[Idea of the proof]
		The action of $\Gdisc$ generates a subgroup $\Gamma$ of $H_2(M,\Z)$
		spanned by fundamental 2-cycles associated with the group action (for
		instance, rectangles in $(q,p)$-coordinates for translations).
		The quantization condition implies that $\int_\Sigma\omega$ lies in a
		discrete subgroup of $\R$, isomorphic to $\Z$ times a fixed generator.
		Define $\hbar$ by requiring that the generator equals $2\pi\hbar$.
		Uniqueness follows from the fact that rescaling $\hbar$ would change the
		quantized values.
	\end{proof}
	
	\begin{remark}[Cohomological interpretation]
		On the quantized cycles themselves $\int_\Sigma\omega\in 2\pi\hbar\,\Z$,
		so $\exp\!\left(i\int_\Sigma\omega/\hbar\right)=1$: the cocycle is
		\emph{trivial} on the lattice, which is exactly why the lattice
		translations commute (Proposition~\ref{prop:braid}). The non-trivial
		structure lives one level up---the continuous Heisenberg group of
		sub-cell displacements carries a non-trivial central extension, a
		class in $H^2$ normalized by $\hbar$---and the integrality condition
		is precisely the statement that this class restricts trivially to the
		discrete subgroup $\Gdisc$.
	\end{remark}
	
	
	% ==================================================
	\section{Geometric canonical commutation relations}
	\label{sec:CCR}
	% ==================================================
	
	We now derive canonical commutation relations from the braiding of
	discrete translations in position and momentum.
	
	\subsection{Discrete translations in phase space}
	
	For simplicity, consider $M=\R^2$ with coordinates $(q,p)$ and
	$\omega=\dd p\wedge\dd q$. Let $a,b>0$ and define translations
	\[
	\tau_q(a):(q,p)\mapsto (q+a,p),\qquad
	\tau_p(b):(q,p)\mapsto (q,p+b).
	\]
	Assume that these generate a $\Z^2$-action contained in $\Gdisc$.
	Fundamental rectangles of size $a\times b$ then have symplectic area
	\[
	\int_{\mathrm{cell}} \omega = \int_0^a \int_0^b \dd p\,\dd q = a b.
	\]
	By Theorem~\ref{thm:hbar} we require $ab=2\pi\hbar$ for a minimal cell.
	
	\subsection{Heisenberg-type relations on states}
	
		Let $\Htw$ be the twin Hilbert space (constructed below). It is
		essential that the Heisenberg phase \emph{cannot} arise if both
		translations are represented as naive coordinate shifts of
		phase-space functions, $\psi(q,p)\mapsto\psi(q-a,p)$ and
		$\psi(q,p)\mapsto\psi(q,p-b)$: such shifts commute identically, since
		both send $\psi(q,p)\mapsto\psi(q-a,p-b)$ irrespective of $a,b$. The
		phase requires the configuration \emph{polarization} (functions of
		$q$ alone), or equivalently the prequantum line-bundle connection
		whose holonomy around a cell is its enclosed symplectic area. In the
		configuration polarization on $L^2(\R_q)$ one sets
		\[
		(T_q(a)\psi)(q) := \psi(q-a) = e^{-iaP/\hbar}\psi,\qquad
		(T_p(b)\psi)(q) := e^{-ibq/\hbar}\psi(q) = e^{-ibQ/\hbar}\psi,
		\]
		so that $T_q(a)$ shifts position while $T_p(b)$ is the momentum
		boost.
	
	Traversing a minimal rectangle in the $(q,p)$-plane yields a phase:
	
		\begin{proposition}[Braiding of translations]\label{prop:braid}
			With the Weyl operators above,
			\begin{equation}\label{eq:braid}
				T_q(a)\, T_p(b) = e^{i ab/\hbar}\, T_p(b)\, T_q(a).
			\end{equation}
			In particular, when the cell area is an integer multiple of the
			minimal area, $ab\in 2\pi\hbar\,\Z$, the phase is
			$e^{iab/\hbar}=1$ and the lattice translations $T_q(a),T_p(b)$
			commute; a non-trivial phase $e^{iab/\hbar}\neq 1$ arises precisely
			for sub-cell (fractional) displacements.
		\end{proposition}
	
	To obtain non-trivial phases in the quantum regime one typically works
	with parameters $(a,b)$ not corresponding to a single minimal cell but
	to fractional displacements relative to the minimal area scale. This
	leads to the Weyl form of the canonical commutation relations.
	
	\begin{theorem}[Geometric canonical commutation]\label{thm:CCR}
		Let $a,b$ be sufficiently small and let $T_q(a)$, $T_p(b)$ be the
		corresponding unitary translations on $\Htw$. Under the quantization
		condition, there exists $\hbar>0$ such that
		\begin{equation}\label{eq:Weyl}
			T_q(a)\, T_p(b) = e^{i ab/\hbar}\, T_p(b)\, T_q(a).
		\end{equation}
		Formally differentiating with respect to $a$ and $b$ at $0$ yields
		densely defined self-adjoint generators $Q,P$ such that
		\[
		[Q,P] = i\hbar \Id
		\]
		on a common invariant core.
	\end{theorem}
	
	\begin{proof}[Sketch of proof]
		The phase factor $e^{iab/\hbar}$ corresponds to the symplectic area
		enclosed by the loop obtained by composing the two translations in
		different orders. The discrete group action quantizes this area in units
		of $2\pi\hbar$. Relation~\eqref{eq:Weyl} is the Weyl form of the
		canonical commutation relations. Expanding $T_q(a)$ and $T_p(b)$ as
		$T_q(a)=e^{-iaP/\hbar}\approx \Id - i a P/\hbar$ and
		$T_p(b)=e^{-ibQ/\hbar}\approx \Id - i b Q/\hbar$ for small $a,b$ and comparing terms
		of order $ab$ yields $[Q,P]=i\hbar\Id$.
	\end{proof}
	
	\begin{remark}
		The construction extends to higher dimensions and to more general
		configuration manifolds, where the role of $q,p$ is played by local
		canonical coordinates and the discrete group action is adapted to the
		underlying geometry.
	\end{remark}
	
	
	% ==================================================
	\section{Twin Hilbert spaces as $L^2$ on orbit space}
	\label{sec:Hilbert}
	% ==================================================
	
	We now construct the Hilbert space of the theory as an $L^2$-space on
	the orbit space.
	
	\subsection{Invariant measures on the quotient}
	
	We assume:
	
	\begin{itemize}[leftmargin=2em]
		\item The action of $\Gdisc$ on $M$ is properly discontinuous and admits
		a fundamental domain $F\subset M$ of finite symplectic volume;
		\item The symplectic volume form $\Omega:=\omega^{\wedge n}$ induces a
		$\Gdisc$-invariant measure on $M$.
	\end{itemize}
	
	The quotient $X=M/\Gdisc$ is then equipped with a natural measure $\mu$
	defined by
	\[
	\int_X f([x])\, \dd\mu([x])
	:= \int_F f(\pi(x))\, \Omega(x)
	\]
	for reasonable test functions $f$.
	
	\begin{proposition}
		The measure $\mu$ is well-defined (independent of the choice of $F$ up
		to null sets) and $\sigma$-finite. If $F$ has finite volume, then
		$\mu(X)<\infty$.
	\end{proposition}
	
	\begin{proof}[Idea]
		The invariance of $\Omega$ under $\Gdisc$ ensures consistency of the
		definition on overlapping translates of $F$. Proper discontinuity
		guarantees that only countably many translates intersect a given compact
		set, leading to $\sigma$-finiteness.
	\end{proof}
	
	
	\subsection{Construction of the twin Hilbert space}
	
	\begin{definition}[Twin Hilbert space]
		The \emph{twin Hilbert space} associated with $(M,\omega,\Gdisc)$ is
		\[
		\Htw := L^2(X,\mu),
		\]
		the space of (equivalence classes of) complex-valued, square-integrable
		functions on the orbit space $X=M/\Gdisc$.
	\end{definition}
	
	\begin{proposition}
		$\Htw$ is a separable Hilbert space with inner product
		\[
		\langle \psi,\varphi\rangle
		:= \int_X \overline{\psi([x])}\,\varphi([x])\,\dd\mu([x]).
		\]
	\end{proposition}
	
	\begin{proof}
		Standard $L^2$ theory on a $\sigma$-finite measure space.
		Separability follows from the existence of a countable dense subset
		(e.g.\ simple functions with rational coefficients).
	\end{proof}
	
	\begin{remark}[Prequantum space vs.\ polarization]\label{rem:polarization}
		The space $L^2(X,\mu)$ built from the \emph{full} phase space
		$X=M/\Gdisc$ is the \emph{prequantum} Hilbert space. As in geometric
		quantization it is reducible---heuristically ``twice too large''---and
		the irreducible physical Hilbert space is obtained by choosing a
		\emph{polarization} (a maximal involutive Lagrangian distribution) and
		keeping only polarized sections. In the configuration polarization the
		physical space is $L^2$ of the configuration quotient, e.g.\
		$L^2(\R_q)$ in Section~\ref{sec:examples}. Throughout we regard
		$L^2(M/\Gdisc)$ as the ambient prequantum space and the physical
		$\Htw$ as its polarized subspace; the Weyl operators of
		Section~\ref{sec:CCR} act in this polarization.
	\end{remark}
	
	\subsection{Representation of translations and observables}
	
	The discrete action and symplectic structure induce natural operators:
	
	\begin{itemize}[leftmargin=2em]
		\item For each $g\in\Gdisc$, a unitary operator $U(g)$ on $\Htw$ defined
		by transport along the orbit; in many examples $U(g)$ is trivial
		at the quotient level but non-trivial when internal phases are
		included;
		\item Multiplication operators $(A\psi)([x]) = a([x])\psi([x])$ for
		suitable measurable functions $a:X\to\R,\C$, corresponding to
		position or other configuration observables;
		\item Differential or pseudo-differential operators representing
		momentum and Hamiltonians, obtained by lifting vector fields on
		$M$ and projecting them to the orbit space.
	\end{itemize}
	
	\begin{theorem}[Twin Hilbert construction]
		Under the assumptions above, the triple $(\Htw,U,\mathcal{A})$, where
		$\mathcal{A}$ is a suitable *-algebra of operators generated by
		observables and translations, provides a representation of the twin
		quantum geometry compatible with the symplectic form and the discrete
		action.
	\end{theorem}
	
	\begin{remark}
		In later works, we refine $\Htw$ by incorporating internal phases and
		cohomological data, leading to projective representations and central
		extensions of $\Gdisc$. This is the natural arena for the full Weyl
		system associated with $(Q,P)$.
	\end{remark}
	
	
	% ==================================================
	\section{Twin time evolution and Schr\"odinger dynamics}
	\label{sec:Schro}
	% ==================================================
	
	We now show how Schr\"odinger-type unitary evolution arises from
	covariant transport of Hamiltonian flows.
	
	\subsection{Hamiltonian flows on $(M,\omega)$}
	
	Let $H_{\mathrm{cl}}:M\to\R$ be a smooth Hamiltonian. The associated
	Hamiltonian vector field $X_H$ is defined by
	\[
	\iota_{X_H}\omega = \dd H_{\mathrm{cl}}.
	\]
	The flow $\Phi_t:M\to M$ generated by $X_H$ preserves $\omega$ and
	hence the symplectic volume form $\Omega$.
	
	We assume that $\Phi_t$ commutes with the action of $\Gdisc$ or, more
	generally, that it is compatible up to phase factors determined by the
	cohomological data.
	
	\subsection{Induced evolution on the orbit space}
	
	The flow $\Phi_t$ descends to a measurable flow
	\[
	\tilde{\Phi}_t : X \to X,\qquad \tilde{\Phi}_t([x]) = [\Phi_t(x)],
	\]
	which preserves the measure $\mu$.
	
	\begin{definition}[Twin time evolution]
		The induced unitary group $(U_t)_{t\in\R}$ on $\Htw$ is defined by
		\[
		(U_t\psi)([x]) := \psi(\tilde{\Phi}_{-t}([x])).
		\]
	\end{definition}
	
	\begin{proposition}
		Each $U_t$ is unitary on $\Htw$, and $(U_t)_{t\in\R}$ is a strongly
		continuous one-parameter group of unitaries.
	\end{proposition}
	
	\begin{proof}
		Unitarity follows from invariance of $\mu$ under $\tilde{\Phi}_t$.
		Strong continuity is a consequence of the continuity of the flow and the
		dominated convergence theorem, for instance on a core of continuous
		compactly supported functions on $X$.
	\end{proof}
	
	\subsection{Generator and Schr\"odinger equation}
	
	By Stone's theorem there exists a self-adjoint operator $\widehat{H}$
	on $\Htw$ such that
	\[
	U_t = e^{-i t \widehat{H}/\hbar}.
	\]
	
	\begin{theorem}[Emergent Schr\"odinger dynamics]
		Let $\psi_t = U_t\psi_0$ for some initial state $\psi_0\in\Htw$. Then
		$\psi_t$ satisfies
		\[
		i\hbar \frac{\partial}{\partial t} \psi_t = \widehat{H}\,\psi_t
		\]
		in the sense of strongly continuous unitary evolution. The operator
		$\widehat{H}$ is the twin-quantized Hamiltonian associated with
		$H_{\mathrm{cl}}$ and the discrete symplectic action.
	\end{theorem}
	
	\begin{proof}
		This is a direct consequence of Stone's theorem. The essential point is
		that the dependence on $\hbar$ has already been fixed by the symplectic
		quantization condition of Section~\ref{sec:hbar}. The generator
		$\widehat{H}$ is determined (up to additive constants) by the classical
		Hamiltonian and the chosen representation of translations and observables
		on $\Htw$.
	\end{proof}

	\begin{remark}[Transport vs.\ Schr\"odinger generator]\label{rem:koopman}
		The unitary $U_t\psi=\psi\circ\tilde\Phi_{-t}$ is the \emph{Koopman}
		operator of the classical flow; its generator on the prequantum space
		is the Liouville operator $\widehat{H}_{\mathrm{pre}}\sim -i\hbar\,X_H$,
		whose spectrum is the classical (Liouville) spectrum, \emph{not} the
		quantum one. Recovering the Schr\"odinger operator---with its discrete
		spectrum and zero-point energy---requires the polarization of
		Remark~\ref{rem:polarization} together with the prequantum connection
		term, i.e.\ the geometric-quantization operator
		$\widehat{H}=-i\hbar\,X_H+\big(H_{\mathrm{cl}}-\theta(X_H)\big)$
		restricted to polarized sections. The statements of this section are
		to be read in that polarized sense.
	\end{remark}
	
	\begin{remark}
		In explicit models (harmonic oscillator, free particle, particle on a
		circle), $\widehat{H}$ recovers the familiar Schr\"odinger operators,
		with $\hbar$ and canonical commutation relations inherited from the twin
		geometry.
	\end{remark}
	
	
	% ==================================================
	\section{Examples and simple models}
	\label{sec:examples}
	% ==================================================
	
	We work out three canonical examples, showing how the symplectic-area quantization of \S\ref{sec:hbar} reproduces the standard spectra; the sketches below are made fully precise in a later
	version.
	
	\subsection{Harmonic oscillator}
	
	Take $M=\R^2$ with coordinates $(q,p)$ and
	\[
	H_{\mathrm{cl}}(q,p) = \frac{p^2}{2m} + \frac{1}{2}m\omega^2 q^2.
	\]
	Choose a rectangular lattice generated by translations
	$(q,p)\mapsto(q+a,p)$ and $(q,p)\mapsto(q,p+b)$ with $ab=2\pi\hbar$.
	Fundamental twin cells have area $2\pi\hbar$. The twin Hilbert space
	$\Htw=L^2(X,\mu)$ carries a representation of $(Q,P)$ satisfying
	$[Q,P]=i\hbar\Id$. The twin Hamiltonian $\widehat{H}$ defined via the
	unitary group $U_t$ reproduces, in the configuration polarization (Remark~\ref{rem:polarization}; cf.\ Remark~\ref{rem:koopman}), the usual spectrum
	$E_n=\hbar\omega(n+\tfrac12)$. Concretely the framework \emph{fixes} this spectrum by area quantization. A
classical orbit $H_{\mathrm{cl}}=E$ is the ellipse $\tfrac{p^2}{2m}+\tfrac12 m\omega^2 q^2=E$, with semi-axes
$\sqrt{2mE}$ (in $p$) and $\sqrt{2E/(m\omega^2)}$ (in $q$), hence enclosed symplectic area
\[
\oint_{H=E} p\,\dd q=\pi\sqrt{2mE}\,\sqrt{\tfrac{2E}{m\omega^2}}=\frac{2\pi E}{\omega}.
\]
Bohr--Sommerfeld quantization in units of the minimal twin cell, with the Maslov index $\tfrac12$ supplied by
the polarization (Remark~\ref{rem:polarization}), reads $\oint p\,\dd q=2\pi\hbar\,(n+\tfrac12)$, i.e.\
$\tfrac{2\pi E}{\omega}=2\pi\hbar(n+\tfrac12)$, giving exactly $E_n=\hbar\omega(n+\tfrac12)$. The zero-point
energy $\tfrac12\hbar\omega$ is the area of half a fundamental cell --- a direct geometric reading of the
paper's central quantity $\Delta\mathcal{A}_{\min}=2\pi\hbar$.
	
	\subsection{Free particle}
	
	For a free particle of mass $m$ on the line, $H_{\mathrm{cl}}(q,p) =
	p^2/(2m)$. The twin structure quantizes the phase space area and
	identifies $\hbar$ as the minimal cell area divided by $2\pi$. The twin
	Hilbert space is $L^2(\R,\dd q)$---the configuration polarization of the prequantum space $L^2(M/\Gdisc)$ (Remark~\ref{rem:polarization})---and the generator of translations in
	$q$ acts as the usual momentum operator $P=-i\hbar\,\partial_q$.
	
	\subsection{Particle on a circle}
	
	Taking $M=T^*S^1\simeq S^1\times\R$ with angle coordinate $\theta$ and
	momentum $p$, the symplectic form is $\omega=\dd p\wedge\dd\theta$.
	Quantization of $\int \omega$ along cycles associated with $\theta
	\mapsto\theta+2\pi$ yields discrete momentum eigenvalues $p_n = n\hbar$.
	This fits naturally in the twin framework, with $\Gdisc$ generated by
	$2\pi$-shifts in $\theta$ and integer shifts in $p$. Explicitly, in the configuration polarization the physical
states are functions $\psi(\theta)$ on $S^1$; a momentum eigenstate $\psi(\theta)=e^{ip\theta/\hbar}$ is
single-valued under $\theta\mapsto\theta+2\pi$ iff $p\cdot 2\pi/\hbar\in 2\pi\Z$, i.e.\ $p_n=n\hbar$ with
$n\in\Z$. Equivalently, the holonomy $\exp\!\big(\tfrac{i}{\hbar}\oint p\,\dd\theta\big)=1$ around the
$\theta$-cycle is the integrality condition of \S\ref{sec:hbar}, and the discrete spectrum $p_n=n\hbar$ is its
solution --- the simplest instance of emergent quantization in the twin framework.
	
	% ==================================================
	\section*{Conclusion}
	
	We have shown how the Planck constant, canonical commutation relations,
	Hilbert space structure and Schr\"odinger evolution can be derived from
	a purely geometric framework based on discrete symplectic group actions
	on phase space. This provides a rigorous and conceptually unified
	foundation for the twin formulation of quantum mechanics. In subsequent
	articles we will extend this construction to a full operator and
	measurement theory, and further to quantum fields, kernels and
	gravity-like effects.
	
\end{document}
